% English translation of capitulos/02-fundamentacao.tex, made by Claude Code (Claude Fable 5.2) on 24 November 2026.
% Pending review by the author. When the Portuguese file changes, update this one (scripts/verificar_traducao.py).
\section{Background and related work}\label{sec:fundamentacao}
\subsection{Language models}
Current LLMs derive, to a large extent, from the Transformer architecture, based on attention
mechanisms~\citep{vaswani2017attention}. These models are pre-trained on large volumes of text and
directly model the probability of text sequences~\citep{liu2021pretrain}. They can therefore
be used in new tasks without changing the parameters, only by changing the input~\citep{sahoo2024survey}.
It is this possibility that separates the focus of this work from approaches based on training or fine-tuning.

\subsection{Prompt engineering}
A prompt is the instruction given to the model to enforce rules, automate processes and obtain certain
qualities in the output~\citep{white2023patterns}. Prompt engineering is the practice of designing these
inputs to obtain the desired behaviour from an already trained model~\citep{sahoo2024survey}.
\citet{white2023patterns} treat prompts as a form of programming and catalogue
reusable patterns, by analogy with software patterns.

\subsection{Prompting techniques}
The techniques below illustrate the kind of object the review intends to classify. The list is not
exhaustive, its presentation here does not indicate which one works best, and the mention does not decide the
eligibility of any family of studies.
\begin{itemize}
  \item \textit{Zero-shot}: only a natural-language instruction, without
  demonstrations~\citep{brown2020fewshot}.
  \item \textit{Few-shot}: the prompt includes a few demonstrations of the task.
  \citet{brown2020fewshot} evaluate this condition, alongside \textit{one-shot} and \textit{zero-shot},
  as ways of measuring in-context learning, so the two terms are not perfect synonyms.
  \item \textit{Chain-of-thought}: in the formulation of \citet{wei2022cot}, the demonstrations
  carry a series of intermediate reasoning steps. Variants without demonstrations will only be
  described if they appear in the corpus.
  \item ReAct: the model interleaves reasoning traces and actions, which allows it to consult external sources or
  environments~\citep{yao2023react}.
  \item \textit{Self-planning}: the model first plans the solution and then writes the code,
  from the plan~\citep{jiang2025selfplanningv5}.
  \item \textit{Self-debugging}: demonstrations teach the model to debug its own program,
  examining execution results and explaining the code~\citep{chen2023selfdebug}.
\end{itemize}
These items are not of the same kind: some describe a way of building the prompt, and others, a
sequence of interaction with the model.
For this reason, the synthesis uses two units. The technique, with the name given by the authors, is the unit of the comparisons.
The mechanism is the unit of the classification: each technique receives one or more labels from a closed dictionary of eleven
mechanisms, which are instruction, context, constraints, examples, decomposition, reasoning, feedback,
sampling, multi-agent, prompt optimization and other. Thus, an interaction sequence such as
\textit{self-debugging} appears as a combination of mechanisms, for example feedback, reasoning and
examples, and not as a category of its own. The classification is multi-label and distinguishes the mechanisms of the complete
package from those the technique adds to the base prompt.

\subsection{Tasks and forms of evaluation}
In this work, generation is the production of code from a description, specification or example.
Translation is the conversion of code between languages or platforms, including migration and transpilation.
Refactoring is the improvement of the structure of the code without a change in behaviour. Debugging and repair
consist of locating and fixing defects. Tests are included when they support these tasks. These
definitions delimit the scope of the review, but do not settle every borderline case:
Section~\ref{sec:metodologia} shows the calibration carried out and the cases still open.

Evaluation usually relies on benchmarks. HumanEval measures the functional correctness of
programs synthesized from docstrings~\citep{chen2021codex}. SWE-bench uses
real GitHub problems, in which the model has to edit a code base to resolve an
issue~\citep{jimenez2024swebench}. Results on these sets do not, by themselves, demonstrate
gains in productivity, security or cost. \citet{sallou2024breaking} discuss
threats to validity in the use of LLMs in software engineering, and part of them guides the quality
assessment planned in this work.

\subsection{Related work}
% AFIRMACAO: A119
\citet{hou2024slr} reviewed 395 studies on LLMs in software engineering, from January 2017 to
January 2024. On prompting, they describe the techniques and count their frequency: \textit{few-shot} in 88 studies,
\textit{zero-shot} in 79 and CoT in 18.
% AFIRMACAO: A120
\citet{fan2023survey} survey the area, rather than conduct a systematic review, and point to hybrid
techniques, which combine traditional software engineering and LLMs, as central to reliable results.
Both cover many software engineering activities, and not only prompting.
% AFIRMACAO: A121
In the opposite direction, \citet{sahoo2024survey} organize 41 prompting techniques by application area,
among them code-specific techniques such as \textit{Program of Thoughts}, SCoT and \textit{Chain of Code}, and
\citet{liu2021pretrain} systematize the prompting paradigm in natural language processing, without
restricting themselves to code tasks.

The search of this work also found closely related secondary studies, recorded as supporting material.
% AFIRMACAO: A122
\citet{yang2025aprsurvey} (S820) classify 66 LLM-based repair systems into fine-tuning,
prompting, procedures and agents and record, for the main result of each system, the benchmark, the
metric, the model and the assumption about fault localization.
% AFIRMACAO: A123
\citet{masthanali2026energy} (S1246) gather 18 studies on the energy efficiency of code
generated by LLMs and conclude that no prompting strategy improves it consistently across configurations.

This work crosses the two perspectives: prompting techniques in the four tasks, with extraction of the type of
counterfactual, the direction of each comparison, the evaluation conditions, the risks and the mitigations. The full
texts of these reviews were read, in the arXiv versions indicated in the references, for the comparison in Section~\ref{sec:comparacao}. Since the
search was not designed to locate reviews, no claim is made that this combination is new.
